% Generated by paper4/manuscript/build.py from paper4/paper/section_partA.md (d5cd4c21240a03dd): edit the Markdown, not this file.
\section{Part A — what does a loop detector measure of congestion?}
Part B found that a sensor buys the measurement at its own segment, and nothing beyond it. Part A asks what that measurement is, for the most common urban sensor. A loop detector measures flow and occupancy; congestion is reconstructed as their ratio, relative to its free-flow value. Three datasets give an independent speed reference: the \textbf{M30 ring of Madrid} (2024), 279 sensors that measure flow, occupancy and speed on the same device; the \textbf{Madrid arterials} (2021), 399 loops matched to probe-speed road segments, a median of 33 m away; and \textbf{Austin}, radar detectors at signalised intersections with volume, occupancy and speed per lane.

Three choices of analysis apply throughout: the congestion threshold is chosen out of sample, every figure has an interval from a block bootstrap, and the effect of the number of lanes is measured as a rank correlation with an interval. The reading rules were fixed knowing the in-sample values of the earlier analysis, so every statement of Part A is descriptive.

\subsection{The loop detects congestion but does not measure its level}
On the M30 ring, on the 4,199,040 cells where the occupancy rule does not impute the value, reconstructed and measured congestion agree poorly in level. The overall rank correlation is +0.249 [+0.208; +0.290]. It changes sign with occupancy, from −0.112 [−0.136; −0.086] at 5–10\% occupancy to +0.463 [+0.284; +0.620] above 40\%, and within the measured classes it is +0.012 [−0.019; +0.044] on free flow, +0.458 [+0.320; +0.617] on slowed traffic and +0.632 [+0.543; +0.718] on congestion.

The same reconstruction separates congested cells from the others almost perfectly: the area under the curve for congestion against the rest is 0.984 [0.979; 0.988], and 0.752 [0.706; 0.804] for slowed or congested against free flow.

\textbf{The threshold transfers between sensors.} A threshold of 0.30 on the reconstructed factor, chosen in sample, gives precision 0.498, recall 0.891 and F1 0.639 on the same cells. Chosen on half of the sensors and reported on the other half, over 5 random splits, it gives an F1 within 0.002 of the in-sample one (0.640 against 0.639):

\begin{table*}[!t]
\centering\footnotesize
\renewcommand{\arraystretch}{1.15}
\caption{The loop detects congestion but does not measure its level}
\begin{tabularx}{\textwidth}{>{\raggedright\arraybackslash\hsize=1.723\hsize}X>{\raggedright\arraybackslash\hsize=0.504\hsize}X>{\raggedright\arraybackslash\hsize=0.924\hsize}X>{\raggedright\arraybackslash\hsize=0.924\hsize}X>{\raggedright\arraybackslash\hsize=0.924\hsize}X}
\toprule
\textbf{Reading} & \textbf{Thresholds chosen (5 splits)} & \textbf{F1} & \textbf{Precision} & \textbf{Recall} \\
\midrule
between sensors, whole year & 0.28–0.29 & 0.640 [0.539; 0.725] & 0.512 [0.397; 0.621] & 0.853 [0.760; 0.926] \\
in time, January–March $\rightarrow$ April–December & 0.23–0.32 & 0.644 [0.539; 0.733] & 0.541 [0.404; 0.666] & 0.810 [0.616; 0.960] \\
\bottomrule
\end{tabularx}
\end{table*}

With the threshold set for 5\% false alarms on the first half, recall on the second half is 0.983 [0.949; 0.999], with a realised false-alarm rate of 0.050 [0.034; 0.083].

\begin{quote}
\textbf{A loop detector says whether a segment is congested, not how congested it is. A threshold chosen on some sensors transfers to others.}
\end{quote}

\subsection{Matching to road segments accounts for about half of the arterial gap}
On the Madrid arterials, loops and probe speeds barely agree: the median over pairs of the per-pair area under the curve is 0.506 [0.504; 0.508] (394 pairs), and the median rank correlation +0.040 [+0.032; +0.048] (399 pairs); on the pooled cells of the pairs, 0.525 [0.517; 0.534] and +0.072 [+0.057; +0.089].

On the M30 ring the reference can be moved away from the loop while keeping regime and technology. Taken from a neighbouring sensor, discrimination falls with distance:

\begin{table*}[!t]
\centering\footnotesize
\renewcommand{\arraystretch}{1.15}
\caption{Matching to road segments accounts for about half of the arterial gap}
\begin{tabularx}{\textwidth}{>{\raggedright\arraybackslash\hsize=2.684\hsize}X>{\raggedright\arraybackslash\hsize=0.331\hsize}X>{\raggedright\arraybackslash\hsize=0.368\hsize}X>{\raggedright\arraybackslash\hsize=0.809\hsize}X>{\raggedright\arraybackslash\hsize=0.809\hsize}X}
\toprule
\textbf{Reference} & \textbf{Sensors} & \textbf{Median distance} & \textbf{Area under the curve} & \textbf{Gini} \\
\midrule
the same sensor & 279 & 0 m & 0.984 [0.979; 0.988] & 0.968 [0.958; 0.976] \\
a neighbour within 0–20 m & 191 & 14 m & 0.789 [0.723; 0.853] & 0.579 [0.447; 0.707] \\
a neighbour within 19–72 m (a posteriori band, kept for the Gini split) & 181 & 31 m & 0.724 [0.660; 0.795] & 0.448 [0.321; 0.591] \\
a neighbour within 20–60 m & 156 & 31 m & 0.764 [0.699; 0.826] & 0.528 [0.398; 0.653] \\
a neighbour within 60–200 m & 165 & 96 m & 0.593 [0.539; 0.655] & 0.187 [0.079; 0.310] \\
\bottomrule
\end{tabularx}
\end{table*}

The Gini scale (twice the area under the curve, minus one) is additive, so the arterial gap can be split. A reference 31 m away removes 0.968 − 0.448 of the Gini on the ring; against the gap down to the arterials (Gini 0.012), that is a share of 0.544 [0.395; 0.676] (0.566 [0.411; 0.705] with the pooled arterial value). The rest belongs jointly to the regime and the technology, which the data cannot separate. The interval includes one half, so the statement is ``undetermined'': matching explains between two fifths and two thirds of the gap.

\subsection{Aggregating lanes hides single-lane episodes}
In Austin, per lane, reconstruction behaves as on the ring: the median within-lane rank correlation is +0.706 [+0.666; +0.764] and the median within-lane area under the curve 0.991 [0.983; 0.998] (interquartile range 0.965–0.998). When the lanes of a section are aggregated, the reconstructed factor must detect an episode on any one of them:

\begin{table*}[!t]
\centering\footnotesize
\renewcommand{\arraystretch}{1.15}
\caption{Aggregating lanes hides single-lane episodes. Sections with $\geq$ 2 lanes: 30 (22/6/2); rank correlation over 29. The difference of prevalence is the share of cells in which at least one lane is congested but the section as a whole is not.}
\begin{tabularx}{\textwidth}{>{\raggedright\arraybackslash\hsize=0.529\hsize}X>{\raggedright\arraybackslash\hsize=0.588\hsize}X>{\raggedright\arraybackslash\hsize=0.647\hsize}X>{\raggedright\arraybackslash\hsize=1.941\hsize}X>{\raggedright\arraybackslash\hsize=0.882\hsize}X>{\raggedright\arraybackslash\hsize=1.412\hsize}X}
\toprule
\textbf{Lanes per section} & \textbf{Sections} & \textbf{Cells} & \textbf{Median area under the curve} & \textbf{Interquartile range} & \textbf{Difference of prevalence (points)} \\
\midrule
2 & 22 & 1,432,457 & 0.990 [0.985; 0.997] & 0.982–0.997 & +0.19 [+0.08; +0.33] \\
3 & 6 & 433,660 & 0.755 [0.461; 0.997] & 0.489–0.995 & +18.62 [+3.87; +33.35] \\
4 & 2 & 154,380 & 0.664 (2 sections, no interval) & 0.533–0.796 & +9.88 (no interval) \\
\bottomrule
\end{tabularx}
\end{table*}

One two-lane section has only 11 aggregated cells with a congested lane, so its area under the curve is not defined (more than 20 cells of each kind are required). Over the other 29 sections, the rank correlation between the number of lanes and the section's area under the curve is −0.441 [−0.714; −0.072]. Aggregation does not change the agreement of volume and occupancy, which is +0.735 [+0.678; +0.829] per lane and +0.739 [+0.582; +0.902] aggregated, over the 8 sections with at least 3 lanes.

\begin{quote}
\textbf{The more lanes are aggregated, the worse the reconstructed factor detects an episode on one of them.}
\end{quote}

For an operator: a loop is a detector of congested episodes, not a meter of their level; its threshold can be set on some sensors and used on others; matching a loop to a road segment some 30 m away costs about half of its apparent disagreement with probe data; and averaging across lanes hides episodes on one lane, whatever the device, so the target definition matters as much as the device.

